\documentclass[10pt,a4paper]{amsart}
\usepackage[margin=19mm]{geometry}
\usepackage{amsmath,amssymb,mathtools}
\usepackage[scr=rsfs,cal=euler]{mathalfa}
\usepackage{xfrac}
\usepackage[unicode,hidelinks,bookmarks=false]{hyperref}
\newcommand{\ie}{i.e.\ }
\newcommand{\cf}{cf.\ }
\newcommand{\resp}{respectively\ }
\newcommand{\Subsection}{\subsection}
\providecommand{\nobreakdash}{\nobreak\hbox{-}}
\setlength{\emergencystretch}{2em}
\multlinegap0pt
\usepackage{bm}
\usepackage{bbm}
\usepackage{dsfont}
\usepackage{xspace}
\usepackage{cancel}
\usepackage{mathtools}
\usepackage{amsthm}
\makeatletter
\@ifundefined{thm}{\newtheorem{thm}{Theorem}}{}
\makeatother
\title[The First Eigenvalue of Embedded Minimal Hypersurfaces in th]{The First Eigenvalue of Embedded Minimal Hypersurfaces in the Unit Sphere I: Yau's Conjecture}
\author{Lingzhong Zeng}
\date{Source: arXiv:2508.06123v1 (CC BY 4.0)}
\begin{document}
\maketitle

\noindent\emph{Source and licence.} The First Eigenvalue of Embedded Minimal Hypersurfaces in the Unit Sphere I: Yau's Conjecture. Version arXiv:2508.06123v1 (CC BY 4.0), distributed under the Creative Commons Attribution 4.0 International licence, as recorded in the arXiv licence metadata for this exact version and shown on the abstract page https://arxiv.org/abs/2508.06123v1. Full rights notice: licenses/arxiv-2508.06123-CC-BY-4.0.txt.\par\smallskip

\noindent\textit{Editorial scope.} Counted unit: the Thm whose source label is \texttt{converse-thm} in the article (source0.tex lines 952--954, source label \texttt{converse-thm}), delivered together with the article's own complete original proof of it (source0.tex lines 957--1032, one \texttt{proof} environment of 76 lines). No mathematics is written, repaired or re-derived: statement and proof are the article's own text line for line. Required context retained uncounted: none. Every definition, display and label of those lines is retained unchanged. Omitted from this package and not redistributed: 1--951, 955--956, 1033--1380. Nothing inside a retained excerpt is omitted or altered: every retained body is delivered exactly as the article's lines are written, so no omission receipt of any kind accompanies this package. Citations: the retained excerpts carry no citation command at all, so no bibliography entry of the article is transcribed and no reference is redistributed. Reference adaptation: none was needed and none was made; every reference inside the delivered unit resolves inside it, so no unresolved reference or citation is produced. Printed slips kept literal: no printed word, symbol, hypothesis or number of the counted statement or of its proof was altered. Layout and class: the article's own class is \texttt{amsart} with packages including mathrsfs, amsmath, amsfonts, amssymb, amsthm, amscd, xy, graphpap, color, eucal, verbatim, hyperref, cite, graphicx; the wrapper is the lane's pinned \texttt{amsart} editorial wrapper at 10pt on A4, which restates the theorem-like environments the retained text needs and adds the article's own macro definitions that the retained text needs (none), with extra wrapper packages (bm, bbm, dsfont, xspace, cancel, mathtools). The delivered numbering is the unit's own. Assets: no retained span contains a figure environment, a graphics-inclusion command, a PDF-page inclusion, a table or a TikZ picture; no image file is copied into this package, the package declares no asset dependency and no image file is redistributed with it. Licence: version arXiv:2508.06123v1 of this article is distributed under the Creative Commons Attribution 4.0 International licence, as recorded in the arXiv licence metadata for this exact version and shown on the abstract page https://arxiv.org/abs/2508.06123v1. A full rights notice is delivered at licenses/arxiv-2508.06123-CC-BY-4.0.txt. Characters: every retained excerpt is pure ASCII, so the delivered engine, XeLaTeX with its default Latin Modern text font, reports no missing character.
\begin{thm}\label{converse-thm}
Let $\Sigma^n \subset \mathbb{S}^{n+1}$ be a closed embedded hypersurface with $\lambda_1(\Delta_\Sigma) = n$. Then $\Sigma$ is minimal.
\end{thm}

 \begin{proof}
To prove that \(\Sigma\) is minimal, we refine the connection between the spectral condition \(\lambda_1(\Delta_\Sigma) = n\) and the extrinsic geometry of \(\Sigma\), with special focus on justifying the Ricci curvature bound involving \(n\):


\textit{Step 1. Eigenfunction and Fundamental Identities.}

Let \(u \in H^1(\Sigma)\) be a first eigenfunction with \(\lambda_1 = n\), then
\begin{equation}\label{eq:eigen}
\Delta_\Sigma u = -n u.
\end{equation}
By the variational characterization of \(\lambda_1\),
\begin{equation}\label{eq:rayleigh}
\int_\Sigma |\nabla_\Sigma u|^2 \, d\mathrm{vol}_\Sigma = n \int_\Sigma u^2 \, d\mathrm{vol}_\Sigma.
\end{equation}


\textit{Step 2. Bochner Formula for the Eigenfunction.}

The Bochner formula relates the Laplacian of the gradient norm to the Hessian and Ricci curvature
\begin{equation}\label{eq:bochner}
\frac{1}{2} \Delta_\Sigma |\nabla_\Sigma u|^2 = |\nabla^2_\Sigma u|^2 + \text{Ric}^\Sigma(\nabla_\Sigma u, \nabla_\Sigma u) + \langle \nabla_\Sigma u, \nabla_\Sigma (\Delta_\Sigma u) \rangle.
\end{equation}
By \eqref{eq:eigen}, \(\nabla_\Sigma (\Delta_\Sigma u) = -n \nabla_\Sigma u\), simplifying \eqref{eq:bochner} to
\begin{equation}\label{eq:bochner_simplified}
\frac{1}{2} \Delta_\Sigma |\nabla_\Sigma u|^2 = |\nabla^2_\Sigma u|^2 + \text{Ric}^\Sigma(\nabla_\Sigma u, \nabla_\Sigma u) - n |\nabla_\Sigma u|^2.
\end{equation}


\textit{Step 3. Ricci Curvature via Gauss Equation.}

For hypersurfaces \(\Sigma \subset \mathbb{S}^{n+1}\), the Gauss equation explicitly connects intrinsic Ricci curvature to extrinsic curvature: second fundamental form \(A\), and the ambient geometry of \(\mathbb{S}^{n+1}\). The ambient sphere \(\mathbb{S}^{n+1}\) has constant sectional curvature \(1\), and hence its Ricci curvature satisfies:
\begin{equation}\label{eq:Ricci cur-sphere}
\text{Ric}^{\mathbb{S}}(X,Y) = n \langle X,Y \rangle \quad \forall X,Y \in T\mathbb{S}^{n+1}.
\end{equation}
For \(\Sigma \subset \mathbb{S}^{n+1}\), the Gauss equation for Ricci curvature is
\begin{equation}\label{eq:gauss}
\text{Ric}^\Sigma(X,Y) = \text{Ric}^{\mathbb{S}}(X,Y) + \langle A(X,X), A(Y,Y) \rangle - |A(X,Y)|^2,
\end{equation}
where \(X,Y \in T\Sigma\). Substituting \eqref{eq:Ricci cur-sphere} gives
\begin{equation}\label{eq:gauss_simplified}
\text{Ric}^\Sigma(X,Y) = n \langle X,Y \rangle + \langle A(X,X), A(Y,Y) \rangle - |A(X,Y)|^2.
\end{equation}


\textit{Step 4. Specializing to \(\nabla_\Sigma u\) and Estimating.}

Set \[X = Y = \frac{\nabla_\Sigma u}{|\nabla_\Sigma u|}\] in \eqref{eq:gauss_simplified}, where \(X\) and \(Y\) are unit vectors wiht \(\nabla_\Sigma u \neq 0\).
Substituting unit vectors \(X = Y\) into \eqref{eq:gauss_simplified} gives
\[
\text{Ric}^\Sigma(X, X) = n   + \left( \langle A(X, X), A(X, X) \rangle - |A(X, X)|^2 \right).
\]
From the equality of the inner product and norm for \(A(X,X)\), for \(\nabla_\Sigma u \neq 0\), we obtain:

\[
\text{Ric}^\Sigma(X, X) = n - |A|^2.
\]
Since \(X = \frac{\nabla_\Sigma u}{|\nabla_\Sigma u|}\), for \(\nabla_\Sigma u \neq 0\), scaling back to the original gradient vector \(\nabla_\Sigma u\) preserves the inequality:
\begin{equation}\label{eq:ricci_bound}
\text{Ric}^\Sigma(\nabla_\Sigma u, \nabla_\Sigma u) = n - |A|^2.
\end{equation}Here, we use the following facts that Ricci curvature is homogeneous of degree 0 in the vector argument.

\textit{Step 5. Integrating the Bochner Formula.}

Integrate \eqref{eq:bochner_simplified} over \(\Sigma\). The left-hand side vanishes by the divergence theorem (since \(\Sigma\) is closed), giving
\[
0 = \int_\Sigma |\nabla^2_\Sigma u|^2 \, d\mathrm{vol}_\Sigma + \int_\Sigma \text{Ric}^\Sigma(\nabla_\Sigma u, \nabla_\Sigma u) \, d\mathrm{vol}_\Sigma - n \int_\Sigma |\nabla_\Sigma u|^2 \, d\mathrm{vol}_\Sigma.
\] Substitute \eqref{eq:ricci_bound} and \eqref{eq:rayleigh} into the equation:
\[
0 = \int_\Sigma |\nabla^2_\Sigma u|^2 \, d\mathrm{vol}_\Sigma + \int_\Sigma (n - |A|^2) |\nabla_\Sigma u|^2 \, d\mathrm{vol}_\Sigma - n^2 \int_\Sigma u^2 \, d\mathrm{vol}_\Sigma.
\]

\textit{Step 6. Hessian Estimate and Conclusion.}

By Cauchy-Schwarz, using \eqref{eq:eigen}, we have \(|\nabla^2_\Sigma u|^2 \geq \frac{1}{n} (\Delta_\Sigma u)^2 = n u^2\). Substituting this and simplifying with \eqref{eq:rayleigh} shows all terms must vanish, forcing \(|A|^2 = 0\) almost everywhere.
Thus, \(A \equiv 0\), which means that the mean curvature \(H = \frac{1}{n} \text{tr}(A) \equiv 0\), proving \(\Sigma\) is minimal.
\end{proof}

\end{document}
